\section{What is a solution to the toy model, formally}\label{app:formal-def-solution}

First we define the space of configurations of cAMP molecules spawned by a single amoeba.
Here the $\star$ means that a particle has not been spawned yet.
We must take only sequences where only the first coordinates (finitely) are different from $\star$.

\begin{definition}
Define the space of the cAMP particles spawned by an amoeba as:
\begin{equation}
\R^d_{camp}:=\boxplus_{g=1}^\N (\{\star\} \sqcup \R^d) := 
\{
	x = (x_g)_{g\in\N}\
	:\
	\exists G\in\N_0.\
	\forall g < G.\
	x_g \in \R^d\
	\text{ and }\
	\forall g\geq G.\
	x_g =\star
\}
\end{equation}
Define the space of configurations for all particles as
\begin{equation}
\Conf:=(\R^d)^N \times\big(\R^d_{camp}\big)^N
\end{equation}
For each configuration of particles $x\in \Conf$ sign its coordinates as 
\begin{equation*}
\begin{aligned}
x &= (x^1, \ldots, x^N, 
	&(x^{1,g})_{g\in\N}, \ldots, 
	&(x^{p,g})_{g\in\N}, \ldots, 
	&(x^{N,g})_{g\in\N})
\\
&= (a^1, \ldots, a^N, 
	&c^{1,1}, c^{1,2}, \ldots, 
	&c^{p,1},\ldots, 
	&c^{N,1}, \ldots)
\end{aligned}
\end{equation*}
where the absence of a coordinate in the list means that such coordinate is set to $\star$.
\end{definition}

This definition is important to have processes and function with domains and codomains which are constant throughout time.
For example:

\begin{definition}\label{ch3:formal-def-drift-diffus}
	Let $S$ be the space of \cadlag increasing functions $f:\R+ \longrightarrow \N_0$ (basically the space of increasing "stairs" with steps of integer height). Then define $\drift: \R_+ \times \Conf \times S^N \longrightarrow \Conf$ and $\diffusion: \Conf \longrightarrow \Conf$ as
\begin{equation}
\left\{
\begin{aligned}
&\drift_p(t,x,(\hmactive_t^p)_p) = \frac{1}{N} \sum _{p = 1}^{N} \sum_{g = 1}^{\hmactive_s^p} \attraction(x^p - x^{p,g})
	&&
\\
&\drift_{p,g}(t,x) = 0
	&& \text{if } \btime^{p,g} \leq t
\\
&\drift_{p,g}(t,x) = \star
	&& \text{otherwise}
\\
&\diffusion_p:= \sigma_A
&
\\
&\diffusion_{p,g} = \sigma_C
&& \text{if } \btime^{p,g} \leq t
\\
&\diffusion_{p,g} = \star
&& \text{otherwise}
\end{aligned}
\right.
\end{equation}
\end{definition}

Now we need to see what does it mean to formally solve a piecewise define SDE with a quasi-brownian motion as ours:

%todo
[[[REWRITE FROM HERE ON]]]


\begin{definition}
A \emph{piecewise SDE with quasi BM} is an expression of the type $E($
\end{definition}



\begin{definition}
A solution to \ref{ch2:model:deterministic-parenting} is given by
\begin{itemize}
\item a complete filtered space $(\Omega, \Flt, (\Flt_t)_t, \Prb)$
\item $N$ indipendent Poisson processes $(\hmactive^p)_{1\leq p\leq N}$ on $\Omega$
\item $(\Flt_t)_t$-adapted processes $\sol,B:\R_+ \longrightarrow \Conf$ such that 
	\begin{itemize}
	\item $(\hmactive^p)_{1\leq p\leq N}$ do not depend on $(\sol,B)$
	\item Under the random measure $\Prb(\quad\cdot\quad \big| (\hmactive^p)_{1\leq p\leq N, t\in\R_+})$ we have that $(\sol,B)$ solves $E(\drift((\hmactive_t^p)_p), \diffusion)$ in the coordinates different from $\star$ piecewise on $I_m=[\btime_m,\btime_{m+1}]$ with the opportune modifications from $B^{p,g}$ to $\quasibm^{p,g}$
	\end{itemize}
\end{itemize}
\end{definition}


\section{Small tricks in stochastics}

\begin{lemma}[Estimates of the moments of a Poisson \rv]
	\label{app:Poisson-moments}
	Let $P \sim \Pois(\mu)$ be a Poisson \rv of parameter $mu$.
	Then there is an increasing function $f^{\uparrow}: \N \longrightarrow \N$ independent from $\mu$ such that 
	\begin{equation*}
		\E{|P|^n} \leq  f^{\uparrow}(n)\: (\mu^n \vee 1)
	\end{equation*}
\end{lemma}

\begin{proof}
	In \cite[Equation (1.3-14)]{handbook-poisson} we have that for each $P \sim \Pois(\mu)$
	
	\begin{equation*}
		\E{
			|P|^n
		}
		=
		\sum_{k=1}^{n}
		\mu^k
		\left\{
		\begin{gathered}
			n
			\\
			k
		\end{gathered}
		\right\}\
	\end{equation*}
	
	where 
	
	\begin{equation*}
		\left\{
		\begin{gathered}
			n
			\\
			k
		\end{gathered}
		\right\}\ := \left(
		\frac{1}{k!}\
		\sum_{i=0}^{k}\
		(-1)^{k-i}\
		\left(
		\begin{gathered}
			k
			\\
			i
		\end{gathered}
		\right)\
		i^n
		\right)
	\end{equation*}
	
	are the Stirling number of the second kind.
	Then it is sufficient to observe that
	
	\begin{equation*}
		\E{
			|P|^n
		}
		=
			\sum_{k=1}^{n}
			\mu^k
			\left\{
			\begin{gathered}
				n
				\\
				k
			\end{gathered}
			\right\}\
		\leq
			\sum_{k=1}^{n}
			(\mu \vee 1)^k
			\left\{
			\begin{gathered}
				n
				\\
				k
			\end{gathered}
			\right\}\
		\leq
			(\mu \vee 1)^n
			\sum_{k=1}^{n}
			\left\{
			\begin{gathered}
				n
				\\
				k
			\end{gathered}
			\right\}\		
		=
			(\mu^n \vee 1)
			\underbrace{
				\sum_{k=1}^{n}
				\left\{
				\begin{gathered}
					n
					\\
					k
				\end{gathered}
				\right\}\
			}_{=:f^{\uparrow}(n)}
	\end{equation*}
	
	to obtain the thesis.
\end{proof}

We have also this useful result from \cite[Equation (1.6-9)]{handbook-poisson}.

\begin{lemma}[The moment generating function of a Poisson \rv]
	\label{app:Poisson-generating-function}
	Let $P \sim \Pois(\mu)$ be a Poisson \rv of parameter $\mu$.
	Then its moment generating function is given by:
	\begin{equation*}
		\E{\exp(sP)} =  \exp[\mu(e^s-1)]
	\end{equation*}
\end{lemma}

The next one is a one liner proof that looks absolutely trivial here. Yet in some very cumbersome calculations, where the \rvs involved are nasty, we might miss that some inequalities happen because of this obvious fact.

\begin{lemma}\label{app:pA-pB-less-than}
	Let $A,B$ be non-negative \rvs and let $a,b \in R_+$.
	Then we have that
	\begin{equation*}
		\PrbOf{A+B \geq a +b} \leq \PrbOf{A \geq a} + \PrbOf{B \geq b}
	\end{equation*}	
\end{lemma}

\begin{proof}
	\begin{align*}
		\PrbOf{A+B \geq a +b} &= 
		\PrbOf{A \geq a, A+B \geq a +b} + \PrbOf{A < a, A+B \geq a +b} 
		=
		\PrbOf{A \geq a, A+B \geq a +b} + \PrbOf{B \geq b, A+B \geq a +b} 
		\leq
		\PrbOf{A \geq a} + \PrbOf{B \geq b}
	\end{align*}	
\end{proof}